\chapter{Conclusiones}

La implementación de la Plataforma de Análisis Estadístico Académico resolvió de manera definitiva la fragmentación y dispersión de la información escolar en la Universidad Politécnica de Texcoco. El sistema transformó por completo la forma en que se manejan los datos dentro de la institución, pasando de un proceso manual en hojas de cálculo que tomaba días a una plataforma rápida y segura que entrega resultados claros en corto tiempo. Esto no solo agiliza el trabajo diario en la oficina académica, sino que brinda a los directivos y coordinadores información real y oportuna para entender mejor el avance de los estudiantes, prevenir la deserción y tomar decisiones que realmente beneficien a la comunidad universitaria.

Para nosotros, desarrollar este proyecto fue el mayor reto de nuestra carrera y una experiencia que nos cambió como futuros ingenieros. Salir de los ejercicios del aula y enfrentarnos a un problema real nos enseñó que programar va mucho más allá de escribir código; implicó escuchar a los usuarios, entender sus problemas del día a día, limpiar información compleja y buscar soluciones prácticas cuando las cosas no salían al primer intento.

Más allá del aprendizaje técnico, este proyecto nos deja la satisfacción de haber creado algo útil y con un impacto positivo para nuestra propia universidad. Nos demostró de lo que somos capaces cuando unimos esfuerzos, dedicación y compromiso. Concluir esta etapa nos llena de orgullo y nos da la seguridad y motivación necesarias para dar el siguiente paso y enfrentar con entusiasmo nuestro camino en el mundo profesional.